\section*{Abstract}

This study examines whether the economic literature's diagnosis of the impact of artificial intelligence (AI) on the labor market has changed with the rise of generative AI. It conducts a systematic review, following the PRISMA 2020 guidelines, of studies published between 2013 and mid-2026 and indexed in Web of Science and Scopus. Out of 3,118 records, 816 studies were included and coded, with language models assisting in screening and extraction, by type of evidence, technology in focus, the mechanisms of the Acemoglu--Restrepo task framework, the sign of the employment effect, and the skill group at risk. The share of studies that locate the risk in high-skill occupations rises from 2.7\% among those published up to 2022 to 13.6\% among those published from 2023 onward, although risk for low-skill workers remains the most frequent diagnosis. The change follows the technological object rather than the calendar: high-skill risk is the modal category among studies of generative AI (48.5\%) and is residual and stable among studies of robotics and automation, which serve as a comparison group. This result supports the task approach, since risk follows the content of the tasks each technology performs. Other changes, such as shorter horizons and a less negative sign, also occur within the robotics literature and reflect the empirical maturing of the field. The sign attributed to the employment effect depends strongly on the type of evidence: occupational exposure studies lean negative, firm-level studies lean positive, and theoretical models and surveys lean ambiguous. The new diagnosis rests, so far, on exposure measures and surveys rather than on observed effects. Limitations regarding search coverage and the use of language models as extractors are discussed.

\textbf{Keywords:} artificial intelligence; generative AI; labor market; employment; task approach; systematic literature review.
\newpage
